The {\tt CWAKE} element combines the physics of a co-located {\tt
WAKE} (longitudinal monopole), one or more {\tt TRWAKE} elements
(transverse dipole and quadrupole), and two optional
position-independent (``constant'') transverse deflecting channels,
into a single element.  All channels are driven from one SDDS file;
particle binning, bunch identification, and the per-bunch loop are
done only once across every enabled channel.  This typically reduces
the per-call overhead by about a factor of two relative to stacking
{\tt WAKE} and one or more {\tt TRWAKE} elements at the same location.

\subsection*{Channels}

Seven wake channels are supported; each is enabled by naming its
column in the input SDDS file.  Channels whose column is left at the
default ({\tt NULL}) are disabled.  At least one channel must be
enabled.

\begin{center}
\begin{tabular}{lllll}
\hline
channel        & drive $(d_x,d_y)$ & probe $(p_x,p_y)$ & kick on     & column param \\\hline
longitudinal   & $(0,0)$           & $(0,0)$           & $\Delta E$  & {\tt WZCOLUMN} \\
$x$ dipole     & $(1,0)$           & $(0,0)$           & $\Delta x'$ & {\tt WXCOLUMN} \\
$y$ dipole     & $(0,1)$           & $(0,0)$           & $\Delta y'$ & {\tt WYCOLUMN} \\
$x$ quadrupole & $(0,0)$           & $(1,0)$           & $\Delta x'$ & {\tt QXCOLUMN} \\
$y$ quadrupole & $(0,0)$           & $(0,1)$           & $\Delta y'$ & {\tt QYCOLUMN} \\
$x$ constant   & $(0,0)$           & $(0,0)$           & $\Delta x'$ & {\tt CXCOLUMN} \\
$y$ constant   & $(0,0)$           & $(0,0)$           & $\Delta y'$ & {\tt CYCOLUMN} \\
\hline
\end{tabular}
\end{center}

The two constant channels apply a uniform transverse deflection that depends only on the bunch current $I(t)$, not on
the centroid position of the bunch or the position of the probe particle.  These model deflecting wakes from asymmetric
structures where the wake voltage is independent of the probe particle's transverse coordinates.

Five of the channels ({longitudinal, $x$/$y$ quadrupole, $x$/$y$ constant}) share the same $I(t)$ driver histogram; this
histogram is built exactly once per bunch regardless of how many of those channels are enabled.  Each transverse-dipole
channel adds one additional moment histogram ($xI(t)$ or $yI(t)$).

\subsection*{Scale factors}

Following the {\tt LRWAKE} convention, every channel has its own
scale parameter on top of the global {\tt FACTOR}:
\begin{itemize}
\item {\tt FACTOR} multiplies every channel.
\item {\tt ZFACTOR} multiplies the longitudinal channel only.
\item {\tt XFACTOR} ({\tt YFACTOR}) multiplies the $x$ ($y$) dipole channel.
\item {\tt QXFACTOR} ({\tt QYFACTOR}) multiplies the $x$ ($y$) quadrupole channel.
\item {\tt CXFACTOR} ({\tt CYFACTOR}) multiplies the $x$ ($y$) constant channel.
\end{itemize}
The effective scale of each channel is {\tt FACTOR} $\times$ the
channel factor $\times$ the pass-ramp factor.

\subsection*{Binning}

Bin grids match those of {\tt WAKE} and {\tt TRWAKE} (centered bins,
$t_{\rm min} = t_{\rm mean} - n_b\,dt/2$ when {\tt N\_BINS} is given,
otherwise autoscale from the beam length with $t_{\rm min} -= dt$).
The wake-function time spacing in the input file fixes $dt$.

\subsection*{Notes}

\begin{enumerate}
\item The beam charge must be specified via a {\tt CHARGE} element
  somewhere upstream in the beamline.  Unlike {\tt WAKE} and {\tt
  TRWAKE}, {\tt CWAKE} has no deprecated charge parameter; omitting
  the {\tt CHARGE} element is an error.
\item The original {\tt WAKE} and {\tt TRWAKE} elements are unchanged
  and remain available.  Existing input decks need no modification.
  However, users are encouraged to switch to \verb|CWAKE|.
\item Bunched-mode operation via {\tt BUNCHED\_BEAM\_MODE=1}; see
  Section~\ref{sect:bunchedBeams}.
\item When only the longitudinal channel is enabled, results are
  bit-identical to a stand-alone {\tt WAKE}.  Likewise for any single
  transverse channel and the corresponding stand-alone {\tt TRWAKE}.
\item When the longitudinal channel is combined with any transverse
  channel, {\tt CWAKE} refreshes $p_z$ from the post-longitudinal
  particle state before applying the transverse kicks, matching the
  behavior of stacked {\tt WAKE} $+$ {\tt TRWAKE} elements where each
  subsequent {\tt TRWAKE} recomputes $p_z$ from the updated particles.
  A small residual ($\sim 10^{-6}$ relative) remains when both
  dipole and quadrupole channels are present, because the baseline
  {\tt TRWAKE}(quadrupole) sees a $p_z$ that has been updated by
  the dipole kick, whereas {\tt CWAKE} uses the same post-longitudinal
  $p_z$ for all transverse channels.
\item The {\tt CXCOLUMN}/{\tt CYCOLUMN} constant deflecting channels
  have no single-element counterpart in {\tt WAKE} or {\tt TRWAKE},
  but they are equivalent to a {\tt TRWAKE} element configured with
  all four {\tt X\_DRIVE\_EXPONENT}, {\tt Y\_DRIVE\_EXPONENT},
  {\tt X\_PROBE\_EXPONENT}, and {\tt Y\_PROBE\_EXPONENT} set to zero.
\item The wake-function file must satisfy the same conventions as
  {\tt WAKE}/{\tt TRWAKE}: time column in seconds, monotonic equal
  spacing, wake samples in $V/C$ (longitudinal and constant
  channels), $V/C/{\rm m}$ (dipole), or $V/C/{\rm m}^2$ (quadrupole).
\end{enumerate}
